\section{重复刺激响应与状态干预}
\label{sec:mechanism}
\subsection{固定连接与局部突触抑制的比较}
首轮实验比较M0与M1在$T=0.5/2$s下的响应。每个条件采用五种输入实现，连续给予12次200ms刺激，随后从相同训练末态独立等待2s或10s，再次给予原刺激。共获得20条序列、280个响应窗口。回读结果确认，同一种子在模型、间隔、重复刺激及恢复探测之间的实际感觉事件逐点一致。

M0的aBN1响应在训练过程中保持恒定：四个实现每次10脉冲，一个实现每次9脉冲，两种间隔结果相同；aDN1和aDN2也未出现计数变化。M1在0.5s间隔下出现明显递减，平均$D$为90.22\%，实现间范围为81.25\%--100\%，五例均符合递减、平台及10s恢复判据。2s间隔下的平均$D$为9.88\%，五例均未达到30\%的递减阈值。

\begin{table}[!htbp]\centering\small\setlength{\tabcolsep}{4pt}
\caption{M0与M1的aBN1响应。除D外，数值为五个输入实现的平均脉冲数，每窗口300ms；$P_2$和$P_{10}$分别为等待2s和10s后的探测响应。}
\label{tab:baseline}
\begin{tabular}{lrrrrrr}\toprule
模型/间隔 & $I$ & $E$ & $L$ & $D$(\%) & $P_2$ & $P_{10}$\\\midrule
\inputrows{data/baseline.tex}\bottomrule
\end{tabular}\end{table}

\newcommand{\BasePanel}[2]{\begin{tikzpicture}\begin{axis}[width=.95\linewidth,height=5.3cm,title={#2},xlabel={刺激序号},ylabel={aBN1脉冲数},xmin=1,xmax=12,ymin=0,ymax=11,xtick={1,4,8,12},grid=major,grid style={black!10},tick label style={font=\small}]
\addplot[black!60,thick,mark=square*] table[x=trial,y=mean]{data/base_static_#1.csv};
\addplot[Blue,thick,mark=*] table[x=trial,y=mean]{data/base_std_#1.csv};
\end{axis}\end{tikzpicture}}
\begin{figure}[!htbp]\centering
\begin{minipage}{.48\linewidth}\BasePanel{500}{$T=0.5$s}\end{minipage}\hfill
\begin{minipage}{.48\linewidth}\BasePanel{2000}{$T=2$s}\end{minipage}
\par\smallskip\small 灰色方块：M0\qquad\textcolor{Blue}{蓝色圆点：M1}
\caption{重复刺激下的aBN1响应。每个点为五个输入实现的均值。两图纵轴相同，未按初次响应归一化；M1的初次响应低于M0。逐实现数值见附录。}
\end{figure}

短间隔训练后，M1的2s探测均值为5.40脉冲，10s为8.20脉冲，后者与初次均值相同。逐实现的10s探测/初次响应比为87.5\%--112.5\%，因此组均值相同并不表示每例恰好恢复到初值。原恢复判据为$R_{10}-L\geq0.20E$且$R_{10}\geq0.80I$。

这些结果表明，在本研究的输入、参数和初态下，原模型的快动力学未产生跨次计数递减，而局部资源耗竭与恢复能够产生递减及自发恢复。该比较尚不涉及其他固定权重网络或其他慢机制。

\subsection{初次响应与下行神经元输出}
M1的初次aBN1均值为8.2脉冲，低于M0的9.8脉冲，下降约16.3\%。资源消耗从首次刺激期间即已发生，说明新增机制不仅影响跨刺激历史，也影响单次刺激的传递效能。

下行读出的初始计数更低。M1的aDN1初次响应为0--2脉冲，均值1.2，其中一例为零。如此低的基线限制了相对递减指标的解释，尤其在末期响应接近零时。因此，aBN1作为主要观测量，aDN1/aDN2用于补充描述下游变化。

\begin{table}[!htbp]\centering\small
\caption{下行神经元的初次与末期响应。末期值先对末三次刺激平均，再跨输入实现平均；单位为每300ms窗口的脉冲数。}
\begin{tabular}{lrr}\toprule
条件 & aDN1首次$\to$末期 & aDN2首次$\to$末期\\\midrule
M0，两间隔 & $2.2\to2.2$ & $3.4\to3.4$\\
M1，0.5s & $1.2\to0$ & $2.2\to0$\\
M1，2s & $1.2\to0.2$ & $2.2\to1.4$\\\bottomrule
\end{tabular}\end{table}

仅比较上述输出，尚难区分资源耗竭、快状态残留和下游通路响应能力的变化。后续实验因而分别干预这几项因素。

\subsection{训练末态的选择性重置}
采用原1701--1705输入，各给予12次短间隔刺激。在最后一个300ms响应窗口结束后继续等待200ms，于$t=6$s保存训练末态。各分支分别执行不干预（sham）、恢复资源、重置快状态或同时恢复两类状态，然后给予相同的A探测刺激。

资源恢复将$x$置1，同时同步事件驱动更新的时间戳。快状态重置将膜电位$v$、驱动$g$、最近发放时刻和不应期标记恢复为模型构建时的值，保持$x$不变；后两项在代码中分别为\path{lastspike}和\path{not_refractory}。同步时间戳用于防止下次事件更新时额外计入重置前的恢复过程。

预设判据要求资源恢复后的响应达到M1初次的至少80\%，快状态重置与sham之差不超过1脉冲；两类状态同时恢复后，全脑事件应与初次刺激一致，比较时按刺激起始对齐事件时间。

\begin{figure}[!htbp]\centering
\begin{tikzpicture}\begin{axis}[width=.87\linewidth,height=5.3cm,xmin=-.4,xmax=4.4,ymin=0,ymax=10,xtick={0,1,2,3,4},xticklabels={初次,训练后,重置快状态,恢复资源,同时恢复},ylabel={aBN1脉冲数},tick label style={font=\small},grid=major,grid style={black!10}]
\addplot[Blue,only marks,mark=*,mark size=3pt,error bars/.cd,y dir=both,y explicit] table[x=x,y=mean,y error minus=minus,y error plus=plus]{data/causal.csv};
\end{axis}\end{tikzpicture}
\caption{训练末态干预后的aBN1响应。点为五个输入实现的均值，误差棒表示最小值与最大值。资源恢复与快状态重置分别进行，均使用相同A探测。}
\label{fig:causal}
\end{figure}

回读90个事件文件后，初次aBN1均值为8.2脉冲，训练后sham为0.4；仅重置快状态后仍为0.4，恢复资源后为8.2。五例同时恢复两类状态后的全脑事件均与初次一致。资源干预恢复的是M1自身的初次响应，而非M0的9.8脉冲基线。

\subsection{旁路刺激与结果解释}
旁路实验直接向aBN1注入10个相隔20ms的事件，刺激总时段为200ms，以避开JO-CE输入。未训练和训练后的aBN1/aDN1/aDN2计数均为10/2/3；每个输入实现的全脑事件记录也一致。因此，在重复刺激使JO-CE至aBN1响应降低后，相关下游通路仍能对直接驱动作出响应。该实验未涉及肌肉或身体反馈，无法评价实际运动疲劳；强旁路响应也不能消除原输入条件下的低计数问题。

先前的$U=0$工程对照与M0完整训练事件等价，表明动态突触的替换实现本身未改变原始传递。结合选择性重置结果，可将M1中的主要可逆历史效应归于资源状态。由于耗竭规则由建模者预先指定，这一结果属于模型内的因果检验，尚不能确定真实JO-CE通路采用了相同机制。神经元适应和抑制性可塑性等替代解释仍需另行比较。
\FloatBarrier
